% Folder 93, batch 06 — archivists' pages 101–120.
% Pass: Opus 5.5 (claude-opus-5-5), 2026-10-10 — first pass, unchecked against the pages by a human.
%
% Contents: his manuscript notes on projective connections on a curve:
% the projective connection c_D attached to a connection D on omega, the
% change-of-coordinate formula (Schwarzian derivative), the equation of
% horizontality 2 lambda' + lambda^2 = mu (author's pp. 2, 2 bis, 3, 4, 5:
% archivists' pp. 102, 103, 104, 105, 101), then the dictionary between
% projective connections, sections of P(P^1(omega)), Lie algebras of type
% sl_2 and local uniformisations (pp. 107-114), and a fair statement of the
% Schwarzian G(t,s) as a cocycle, its invariance under PGL(2), and the
% extension to the relative case X/S (pp. 115-120).
% Every page of the batch is in his hand; none was skipped.
\documentclass[11pt,a4paper]{article}
\input{../preamble/grothendieck}

\folder{93}
\batch{6}
\pages{101}{120}
\foldertitle{Connexions de Gauss-Manin et équations de Picard-Fuchs. Opération de Cartier : copies de tapuscrit annoté (s.d.), tiré à part (1981), notes manuscrites (s.d.).}
\dating{[1972]-1981}
\watermark{Édition de démonstration}

\begin{document}

\page{101}
\note{p. 5 de l'auteur (« 5 » cerclé en haut à droite). Dans sa numérotation, cette feuille fait suite à la page 105 du fonds (sa p. 4) ; ses pp. 2, 2 bis, 3 et 4 sont les pages 102 à 105. Sa p. 1 n'est pas dans ce lot : l'argument commence avant lui.}
\[
\frac{dt}{ds}=\frac{1}{ds/dt}=\frac{1}{s'} \qquad ds=s'\,dt
\]
\[
\text{\struck{$d\frac{1}{s}$}}\quad d\,\frac{dt}{ds}=-\frac{s''}{s'^2}\,dt
\qquad ds\,d\,\frac{dt}{ds}=-\frac{s''}{s'}\,dt^2
\]
\[
c_{D_s}=c_{D_t}-2\Bigl(\frac{s''}{s'}\Bigr)'dt^2+\Bigl(\frac{s''}{s'}\Bigr)^2dt^2
\]
\[
\boxed{c_{D_s}=c_{D_t}+\Bigl[-2\Bigl(\frac{s''}{s'}\Bigr)'+\Bigl(\frac{s''}{s'}\Bigr)^2\Bigr]dt^2}
\]
\[
\Bigl[\;s'=\frac{ds}{dt},\quad s''=\frac{ds'}{dt},\quad \Bigl(\frac{s''}{s'}\Bigr)'=\frac{d}{dt}\Bigl(\frac{s''}{s'}\Bigr)\Bigr]
\]
\[
[\;\;]=-2\,\frac{s'''s'-s''^2}{s'^2}+\frac{s''^2}{s'^2}=\frac{3s''^2-2s's'''}{s'^2}
\]
\[
\boxed{c_{D_s}=c_{D_t}+\frac{3s''^2-2s's'''}{s'^2}\,dt^2}
\]

\uncertain{Objection} : $c_{D_s}=c_{D_t}$ ssi \emph{ou bien} $s''\equiv0$ i.e.\ $s$ fonction \ill{} \uncertain{affine}, ou
\[
3s''^2=2s's'''\quad\text{i.e.}\quad 2\,\frac{s'''}{s''}=3\,\frac{s''}{s'}
\]
\note{après le $3$ du second membre, un $/2$ biffé.}
\[
\text{i.e.}\quad (2\log s'')'=(3\log s')'
\]
i.e.\ (\uncertain{cas} $\neq0$)
\[
2\log s''=3\log s'+\mathrm{c}^{\mathrm{te}},\qquad \log\frac{s''^2}{s'^3}=\mathrm{c}^{\mathrm{te}}
\]
\[
\frac{s''^2}{s'^3}=\mathrm{c}^{\mathrm{te}}\quad\text{i.e.}\quad s''^2=Cs'^3,\qquad s''=Cs'^{3/2}
\]
\marginal{entouré, dans la marge gauche : « $s$ fonction homographique de $t$ »}

\[
s=\frac{at+b}{ct+d}=u_M(t)\qquad M\ \text{la matrice}\ \begin{pmatrix}a&b\\c&d\end{pmatrix}
\]
\[
D_{u_M(t)}=D_t
\]
\ill{} \uncertain{peut} \ill{} \ill{} \ill{}.

$3$ types d'homographies : $t\mapsto t+a$, $t\mapsto bt$, $t\mapsto\frac{1}{t}$.
\[
s=\frac1t,\quad s'=-\frac{1}{t^2},\quad s''=+\frac{2}{t^3},\quad s'''=-\frac{6}{t^4}
\]
\[
3s''^2-2s's'''=\frac{12}{t^6}-\frac{12}{t^6}=0
\]

\page{102}
\note{p. 2 de l'auteur (« 2 » cerclé).}
\[
\underline{\omega}+\mathcal{O}_X\xi_0+\underline{\omega}^{-1}v_{s_0}
\]
\[
s=s_0+a_1=(\exp\theta_{-a_1})s_0=s_0+[-a_1,s_0]=\dots
\]
\[
\Bigl[\;v_s=(\exp\theta_{-a_1}\otimes\mathrm{id}_{\underline{\omega}})v_{s_0}
=v_{s_0}+\underbrace{[-a_1,v_{s_0}]}_{-2a_1 s_0}+\tfrac12\bigl[-a_1,\underbrace{[-a_1,v_{s_0}]}_{-2a_1}\bigr]
\]
\[
v_s=\underbrace{-a_1^2}_{2}-\underbrace{2a_1}_{1}+\underbrace{1}_{0}\;\Bigr]
\]
\note{devant $-a_1^2$, un terme biffé (« $4a$ » ?) ; sous l'accolade de droite, « $-\frac12 a_1^2$ », lecture incertaine. Les indices $2,1,0$ sous les termes sont les degrés.}

Je fais calculer, pour avoir l'équation d'horizontalité sur $s$ (ou sur $a_1$) relativement à $c$, la \struck{coeff.\ de} \add{composante} dans $\underline{\omega}^2$ (premier \uncertain{coeff.}) de
\[
(\exp(\theta_{a_1})\otimes\mathrm{id}_{\underline{\omega}})(c(s))
\]
et l'égaler à $0$.

\struck{Ox. \ill{} $c(s)=c(s_0)+c(a_1)$ \ill{} \ill{} pour laquelle un splitting \ill{} horizontal}
\note{passage biffé en zigzag, à peine lisible.}
\[
c(s_0)=\text{\struck{\ill{}}}\;\varpi_2+v_{s_0}
\]
\[
(\exp\theta_{a_1}\otimes\mathrm{id}_{\underline{\omega}})(c(s))
=(\exp\theta_{a_1}\otimes\mathrm{id}_{\underline{\omega}})(c(a_1))+\dots
\]
\[
c(s_0)+c(a_1)=-\varpi_2+c(a_1)+v_{s_0}
\]
\note{la dernière ligne est entre crochets ; le signe devant $\varpi_2$ est surchargé.}

\ill{} \ill{} composantes \ill{} \ill{} $\underline{\omega}^2$, et \ill{}
\[
c(a_1)=\text{\struck{\ill{}}}\;\underbrace{\mathcal{D}_1(a_1)}_{2}-\underbrace{2a_1}_{1}\;\text{\struck{$v_{s_0}$}}\;(\text{\struck{\ill{}}}\dots
\]
sous $\mathcal{D}_1(a_1)$ : « op.\ diff.\ d'ordre $1$, $\underline{\omega}\to\underline{\omega}^2$ » ; à droite : « \ill{} \ill{} splitting \ill{} convenable \ill{} \ill{} \ill{} \ill{} $\delta$ de $\omega_1$ \ill{} $\mathcal{D}_1(\delta)=0$, i.e.\ $\mathcal{D}_1(\lambda\delta)=\delta\,d\lambda$ ».

Et on obtient
\[
(\exp(\theta_{a_1})\otimes\mathrm{id}_{\underline{\omega}})(c(s))\|=\varpi_2+\underbrace{\mathcal{D}_1(a_1)}_{2}-\underbrace{2a_1}_{1}-\underbrace{2[a_1,a_1]}_{+2a_1a_1}
\]
\[
\text{\struck{$+2a$}}\;\underbrace{-a_1^2}_{2}+2a_1+1
\]

L'équation est donc
\[
\boxed{\mathcal{D}_1(a_1)+a_1^2=\varpi_2}
\]
\note{devant $\varpi_2$, un signe surchargé (« $=$ » sur « $-$ ») ; après $\varpi_2$, un indice biffé.}
\marginal{à droite de l'encadré : « équation d'horizontalité de la section de $P$ \uncertain{différente} de $\sigma$, définie par la $1$-forme $a_1$ »}
\marginal{à gauche : « facteur $2$ qui manque — erreur de calcul ? »}

Si on choisit convenablement une base de $\underline{\omega}$ (p.\ ex.\ $\delta=dt$) on peut supposer $\mathcal{D}_1(\delta)=0$ ; dans l'équation en termes de $a_1=\lambda\delta$ est
\[
(d\lambda)\delta+\lambda^2\delta^2=0
\]
en coordonnée, si $\delta=dt$ i.e.\ $a_1=\lambda\,dt$ :
\[
\boxed{\frac{d\lambda}{dt}+\lambda^2=A}
\]
\note{« T.S.V.P. » en bas à droite.}

\page{103}
\note{p. « 2 bis » de l'auteur ; le reste de la feuille est blanc.}
Conclusion. — changer de connexion projective revient simplement \struck{à} \add{à} changer le second membre de l'équation.

\page{104}
\note{p. 3 de l'auteur (« 3 » cerclé).}
\[
\text{Sections de }P-\sigma(X)=\text{connexions sur }\underline{\omega}
\]
$=$ op[érateurs] [diff.\ d'ordre $1$] $\underline{\omega}\to\underline{\omega}^2$ satisfaisant $D(f\omega)=fD\omega+\omega\,df$ ; forment torseur sous $\underline{\omega}$.

À toute telle section correspond une connexion projective particulière $c_D$ — l'unique connexion projective pour laquelle cette section est horizontale. \add{Une section de $P-\sigma(X)$, i.e.\ une section du torseur des connexions $D_1:\underline{\omega}\to\underline{\omega}^2$, i.e.\ un splitting de $\mathcal{L}$.} Si on se fixe une section (p.\ ex.\ si on se donne une connexion sur $\underline{\omega}$, \ill{} \ill{} \ill{}), soit $D_0$, \struck{\ill{}} \struck{lien : \ill{}}
\[
\boxed{c_{D_0+\omega_1}=c_{D_0}+2D_0(\omega_1)+\omega_1^2}
\]
\note{sous $\omega_1$ : « $\hookrightarrow$ \uncertain{conn[exion]} » ; devant le second membre, un essai biffé illisible.}

[Vérification]
\[
c_{D_0+\omega_1+\varpi_1}=c_{D_0+\omega_1}+2(D_0+\omega_1)(\varpi_1)+\varpi_1^2
\]
\[
=\underbrace{c_{D_0}+2D_0\omega_1+\omega_1^2}+2D_0(\varpi_1)+2\omega_1\varpi_1+\varpi_1^2
\]
\[
=c_{D_0}+2D_0(\omega_1+\varpi_1)+\omega_1^2+\varpi_1^2+2\omega_1\varpi_1
=c_{D_0}+2D_0(\omega_1+\varpi_1)+(\omega_1+\varpi_1)^2
\]
« OK ».

Équation d'horizontalité de $D_0+\omega_1$ pour une connexion proj.\ $c_{D_0}+\varpi_2$ :
\[
c_{D_0}+\varpi_2=c_{D_0+\omega_1}\quad\text{i.e.}\quad
\boxed{2D_0(\omega_1)+\omega_1^2=\varpi_2}
\]

Considérons la connexion
\[
D_t(\lambda\,dt)=\frac{d\lambda}{dt}\,dt^2\quad\text{i.e.}\quad D_t(\omega)=\Bigl(\frac{d}{dt}\Bigl(\frac{\omega}{dt}\Bigr)\Bigr)dt^2
\]
On a \struck{$D_0(f\lambda\,dt)$} $=\bigl[d\bigl(\frac{\omega}{dt}\bigr)\bigr]dt$.

On a bien \struck{$D_t(f\omega)=fD_t(\omega)+$}

\note{Dans la moitié inférieure droite, écrit tête-bêche et traversé par trois traits, un calcul séparé de l'action d'un élément $a=\alpha_1+\alpha_{-1}$ sur une base $\omega_1,\lambda,\omega_{-1}$ ; coupé par le bord, il n'est que partiellement lisible :}
\[
\begin{array}{l}
\alpha_1+\alpha_{-1}=a\\
\rho_a\,\omega_1=\text{\struck{$\omega_1\alpha_{-1}$}}\;\dots\;+\lambda\alpha_{-1}\;\dots\\
\rho_a\,\lambda=\lambda\alpha_1+\lambda\alpha_{-1}\qquad
\rho_a\,\omega_{-1}=\alpha_1\omega_{-1}\;\dots\\
(\rho_a)^2\omega_1=(\alpha_{-1}\alpha_1)\omega_1+(\alpha_{-1})^2\omega_2\;\dots\\
(\rho_a)^2\lambda=2\alpha_1\alpha_{-1}\lambda\\
(\rho_a)^2\omega_{-1}=(\alpha_1^2)\omega_{-1}+(\alpha_1\alpha_{-1})\omega_{-1}\\
(\rho_a)^3\omega_1=2\alpha_{-1}^2\alpha_1\omega_1\\
(\rho_a)^3\lambda=2\alpha_1^2\alpha_{-1}\lambda+2\alpha_1\alpha_{-1}^2\lambda
\end{array}
\]
\note{sous ce tableau, la mention « $\underline{\omega}\ \ \underline{\omega}^0\ \ \underline{\omega}^{-1}$ » ; transcription des indices incertaine.}

\page{105}
\note{p. 4 de l'auteur (chiffre cerclé, lu « 4 »).}
\[
D_t(f\omega)=(\text{\struck{\ill{}}}f)D_t\omega+f\omega\,df
\]
\note{il veut sans doute $\omega\,df$ ; le $f$ devant $\omega\,df$ est d'une écriture rapide.}
\[
\left[\;
\begin{aligned}
&d\Bigl(\frac{f\omega}{dt}\Bigr)dt &&\qquad f\Bigl(d\frac{\omega}{dt}\Bigr)dt+\omega\,df\\
&\Bigl(f\,d\frac{\omega}{dt}+df\cdot\frac{\omega}{dt}\Bigr)dt\\
&f\,dt\,d\frac{\omega}{dt}+df\,\omega
\end{aligned}
\;\right]
\]
\note{les lignes de gauche sont reliées par des « $\|$ » verticaux (égalités successives).}

L'équation d'horizontalité pour une \struck{section} connexion $D_t+\lambda\,dt$, relativement à une connexion projective $c_{(D_t)}+\varpi_2$, $\varpi_2=\mu\,dt^2$ :
\[
\text{\struck{$2\lambda'+\lambda^2$}}\qquad
\boxed{2\frac{d\lambda}{dt}+\lambda^2=\mu}
\]

Comparons les connexions projectives $c_{D_t}$ et $c_{D_s}$ quand $ds$ et $dt$ sont bases de $\underline{\omega}$.
\[
\text{\struck{$D_s(\omega)=ds\,d\frac{\omega}{ds}=\dots=\varpi_1\frac{dt}{dt}$}}
\qquad
\text{\struck{$D_t(\omega)=\omega\,d\frac{\omega}{dt}$}}
\]
\[
\text{\struck{$D_s(\omega)-D_t(\omega)=\omega\,d\bigl(\frac{\omega}{ds}-\frac{\omega}{dt}\bigr)$}}
\]
\[
D_s\omega=ds\,d\frac{\omega}{ds}=\frac{ds}{dt}\,dt\,d\Bigl(\frac{\omega}{dt}\,\frac{dt}{ds}\Bigr)
=dt\Bigl[\frac{ds}{dt}\frac{dt}{ds}\,d\frac{\omega}{dt}+\frac{ds}{dt}\,\frac{\omega}{dt}\,d\frac{dt}{ds}\Bigr]
\]
\[
=\underbrace{dt\,d\frac{\omega}{dt}}_{D_t(\omega)}+\text{\struck{\ill{}}}\;\frac{ds}{dt}\,d\frac{dt}{ds}
\]
\[
D_s\omega=D_t\omega+\omega\,\frac{ds}{dt}\,\frac{d}{dt}\Bigl(\frac{dt}{ds}\Bigr)\qquad\text{donc}\qquad
s'\,d\frac{1}{s'}=-\frac{ds'}{s'}
\]
\[
\text{\struck{$c_{D_s\omega}=c_{D_t\omega}+2D$}}\qquad
D_s=D_t+\frac{ds\,d\frac{dt}{ds}}{dt}\quad\text{donc}
\]
\[
c_{D_s}=c_{D_t}+2\,dt\,d\Bigl(\frac{ds\,d\frac{dt}{ds}}{dt^2}\Bigr)+\Bigl(\frac{ds\,d\frac{dt}{ds}}{dt}\Bigr)^2
\]
\note{la suite est p. 101 du fonds (sa p. 5).}

\page{106}
\[
c_{D_0+\omega_1}=c_{D_0}+2D_0(\omega_1)+\omega_1^2
\]

\emph{N.B.} $\omega_1\mapsto 2D_0(\omega_1)+\omega_1^2$ définit un morphisme de faisceaux (ni linéaire, ni « \uncertain{semi-linéaire} »…) $\underline{\omega}\to\underline{\omega}^{\otimes2}$ ; si on remplace $D_0$ par $D_0'=D_0+\varpi_1$, $\varpi_1\in\Gamma(\underline{\omega})$, cela équivaut à $\omega_1\mapsto\omega_1-\varpi_1$.
\note{la cible de $\omega_1\mapsto$ est écrite « $\omega_1\overline{\varpi}$ », avec un trait sur le signe ; lecture incertaine.}

Considérons maintenant une section $\sigma'$ de $P$ sur $X$, donnée (\add{en termes de $D_0$}) par une $\omega_1$ \emph{rationnelle} (pouvant avoir des pôles) — on se demande s'il existe une connexion projective rendant $\sigma'$ plate — laquelle sera unique, on devra avoir \add{$c_{\omega_1}$} $=$ \uncertain{égal} \ill{} $2D_0(\omega_1)+\omega_1^2$. Donc une telle conn.\ projective existe ssi
\[
2D_0(\omega_1)+\omega_1^2\ \text{reste régulière}
\]
aux pôles de $\omega_1$. Si $D_0=d_t$, et $\omega_1=\lambda\,dt$, cela signifie que
\[
2\frac{d\lambda}{dt}+\lambda^2\ \text{reste}
\]
régulière aux pôles de $\lambda$ — cela implique que la \uncertain{partie polaire} de $\lambda$ en un pôle ($t=0$ disons…) est donnée par
\[
\lambda=\frac{2}{t}+\text{termes d'ordre}\geqslant1
\]
(le terme constant devant être nul !)
\note{dans $\frac{2}{t}$, le $2$ est surchargé (sur un $1$) ; au-dessus du « $+$ », un signe biffé.}

\page{107}
Sections horizontales de $P(\underline{P}^1_{X/S}(\omega))$ pour $c_{D_t}$ — i.e.\ connexions $D_t+\omega_1$ ($\omega_1=\lambda\,dt$) horizontales pour $c_{D_t}$, i.e.\ telles que
\[
2\lambda'+\lambda^2=0
\]
Connaissant $D_t$ $\to$ l'uniformisation de la connexion projective canonique sur $\mathbb{P}^1_S$ — \ill{} pour l'uniformisation \emph{\uncertain{globale}} $f:X\to E^1_S\subset\mathbb{P}^1_S$ — \ill{} existe bien qu'il n'y ait pas de connexion \uncertain{partout} définie de $\Omega^1_{\mathbb{P}^1_S/S}$ !) en termes des solutions, ce qui revient \uncertain{aussi} des sections horizontales de $\mathbb{P}_{\mathbb{P}^1_S/S}$ —
\note{reste de la feuille blanc.}

\page{108}
\note{Une liste de descriptions équivalentes d'une même structure sur $X/S$ ; les lettres sont cerclées par l'auteur.}

(a) $(P,\sigma,c)$ : $P$ fibré en droites projectives sur $X$, $c$ connexion sur $P$, \uncertain{rel.} $/S$, $\sigma$ section partout \uncertain{transversale} au \uncertain{stationnaire} \ill{}.

(a') triples $(P,\sigma,\varphi)$, $(P,\sigma)$ comme dessus, $\varphi$ iso \struck{\ill{}} $\underline{\mathcal{O}}_{(\widehat{P},\sigma)}\simeq P^\infty_{X/S}$ d'Alg.\ augm.\ tel que la connexion transportée sur $\mathcal{O}_{(\widehat{P},\sigma)}$ \struck{pour} se globalise à $P$.

(b) Algèbres de Lie de type $\mathfrak{sl}(2)$, quasi-scindée $\mathcal{L}_{-1}=\Omega\subset\mathcal{L}\subset\mathcal{U}_X$ ($\underline{\text{NB}}$ $\mathcal{L}/\Omega\simeq\underline{\mathcal{O}}_X$), $\mathcal{U}\simeq\mathrm{Sym}^2(\mathcal{L})\otimes\Omega^{-1}$, munie de connexion $c:\mathcal{U}\to\mathcal{U}\otimes\underline{\omega}$ (nécessairement $c(\Omega)\subset\mathcal{L}\otimes\underline{\omega}$) telle que le composé $\Omega\xrightarrow{c_\Omega}\mathcal{L}\otimes\omega\to\omega$ soit un iso. ($\underline{\text{NB}}$ Il suffit de se donner $c$ sur $\mathcal{L}$ \struck{\ill{}}, ou $\mathcal{L}\xrightarrow{c_{\mathcal{L}}}\mathrm{Sym}\,\mathcal{L}\otimes\Omega^{-1}\otimes\omega\simeq\mathrm{Sym}^2\mathcal{L}$ — « compatible au crochet » i.e.\ …).

(b') Sur l'Algèbre de Lie quasi-scindée $\mathcal{U}_X$ $\simeq\mathrm{Sym}^2(P^1_{X/S}(\underline{\omega})\otimes\underline{\omega}^{-1})\otimes\underline{\omega}^{-1}$ $\simeq\mathrm{Sym}^2(P^1_{X/S}(\underline{\omega}))\otimes\underline{\omega}^{-3}$, \struck{une connexion} (avec $\mathcal{L}\subset P^1_{X/S}(\underline{\omega})\otimes\underline{\omega}^{-1}$, $\Omega\simeq\underline{\omega}^2\otimes\underline{\omega}^{-1}\simeq\underline{\omega}$), une connexion $c$ d'\emph{Algèbre de Lie} \struck{\ill{}} $c:\mathcal{U}_X\to\mathcal{U}_X\otimes\underline{\omega}$ qui, sur $\underline{\omega}$, soit $d^1_{X/S,\underline{\omega}}:\underline{\omega}\to P^1_{X/S}(\underline{\omega})$.

(c) Posant $\mathbb{P}_X=\mathbb{P}(P^1_{X/S}(\underline{\omega}))$, muni de $\sigma_X$, et dlt.\ \struck{iso} canon.\ $\underline{\mathcal{O}}_{\mathbb{P}_X}/\mathcal{J}^3\overset{\varphi_2}{\simeq}P^2_{X/S}$, un iso
\[
\underline{\mathcal{O}}_{\mathbb{P}_X}/\mathcal{J}^4\overset{\varphi_3}{\simeq}P^3_{X/S}
\]
qui le prolonge.

(d) Une équa.\ diff.\ \add{du $3^{\text{e}}$ ordre rel.\ $X$} \struck{rel.\ $X$} sur $\mathbb{P}^1_S\times_S X\simeq\mathbb{P}^1_X$ (à ne pas confondre avec $\mathbb{P}_X$ !), invariant par $\mathrm{GP}(1)_X$ opérant sur $\mathbb{P}^1_X$, — équa.\ par \uncertain{rapport} aux jets d'ordre $3$ de sections \uncertain{\ill{} isomorphismes} $X\to\mathbb{P}^1_S$ qui \ill{} des iso les \uncertain{\ill{}} d'ordre $3$
\note{la phrase s'arrête au bas de la feuille.}

\page{109}
(e) (Si $S=\operatorname{Spec}(\mathbb{C})$) \add{$X^{\mathrm{an}}$ compact \ill{}} Faisceau d'applications analytiques de $X^{\mathrm{an}}$ dans $\mathbb{P}^1(\mathbb{C})$, qui soit un torseur (à gauche) \add{le faisceau envisagé} sous $\mathrm{GP}(1,\mathbb{C})_X$, et formé d'iso.\ locaux analytiques.

On déduit ces structures équivalentes \struck{connexions} \add{\uncertain{différentes}} \uncertain{(à base de} $X/S$) en termes de \uncertain{d)} à e) en prenant les faisceaux des germes de solutions de l'équa.\ diff., on trouve que \uncertain{ce} sont des solutions locales \uncertain{nécessairement} (la \uncertain{forme} unique bien sûr) par op.\ de $\mathrm{GP}(1,\mathbb{C})$. On passe de e) à a) en \uncertain{considérant} le torseur à gauche \add{\uncertain{correspondant}} (sur $X^{\mathrm{an}}$, du groupe $\mathrm{GP}(1,\mathbb{C})=\mathrm{GP}(1)_{\mathbb{C}}^{\mathrm{an}}$, déduit du torseur \struck{\ill{}} \add{de (e)}, et transformant en torseur de droites par opération symplectique, et prenant le fibré projectif associé. On trouve fibré en droites projectives $P^{\mathrm{an}}$ sur $X^{\mathrm{an}}$ avec connexion analytique — est une classe d'équivalence des couples $(f,y)$, $f$ \uncertain{germe} de section de $\mathcal{U}$ en $x$, $y$ pt de $\mathbb{P}^1(\mathbb{C})$, avec
\[
\text{\struck{$(gf,y)\simeq(f,g^{-1}y)$}}\qquad (f,y)\simeq(\gamma\circ f,\gamma(y))\ \text{pour}\ \forall\,\gamma\in\mathrm{GP}(1,\mathbb{C}).
\]
\note{la ligne biffée et la ligne retenue diffèrent par le côté où opère le groupe.}
Mais dans les couples $(f,\delta f(x))$ ($f\in\mathcal{U}_x$) sont tous conjugués, d'où on trouve une section canonique de $\mathbb{P}^{\mathrm{an}}_x$, d'où une section de $P^{\mathrm{an}}$ sur $X^{\mathrm{an}}$, d'où — \uncertain{un \ill{}} $\forall$ de suite — qui elle est analytique et nulle part stationnaire \uncertain{puisque} les $f\in\mathcal{U}_x$ \uncertain{sont} écrits locaux \struck{\ill{}} \add{\ill{}}.
\[
P|\mathcal{U}\simeq\mathbb{P}^1(\mathbb{C})\times\mathcal{U},\quad\text{la section s'identifie à}
\]
l'application $f:\mathcal{U}\to\mathbb{P}^1(\mathbb{C})$ au voisinage de $x$. On termine par GAGA.

On passe de a) à e) en utilisant la connexion $c^{\mathrm{an}}$ de $\mathbb{P}^{\mathrm{an}}$ sur $X^{\mathrm{an}}$ pour identifier \struck{\ill{}} $X^{\mathrm{an}}|\mathcal{U}$, ($\mathcal{U}$ \uncertain{près} voisinage d'un $x\in X^{\mathrm{an}}$)
\note{la phrase se poursuit p. 110.}

\page{110}
\note{suite de la phrase de la p. 109.}
à un produit $\mathcal{U}^{\mathrm{an}}\times\mathbb{P}^{\mathrm{an}}_x$, donc \struck{\ill{}} s'identifie à une application $\mathcal{U}^{\mathrm{an}}\to\mathbb{P}^{1\,\mathrm{an}}_X$ dans les composés avec les iso.\ $\mathbb{P}^1_X\to\mathbb{P}^1(\mathbb{C})$ donnent les germes d'applications \struck{\ill{}} \add{\uncertain{locales}} \uncertain{du} faisceau $\mathcal{U}$.

Quelle est l'interprétation des sections horizontales \add{locales} (de $P^{\mathrm{an}}$) dans la description (e) ? Pour \uncertain{le} fibré principal (torseur) $\mathcal{T}$ sur $X$ [\uncertain{donnant}
\[
P^{\mathrm{an}}=\underline{\mathrm{Isom}}(\mathbb{P}^1(\mathbb{C})_{X^{\mathrm{an}}},P)\ ]
\]
$P^{\mathrm{an}}$ sur $X^{\mathrm{an}}$, il faut prendre les sections horizontales de $P^{\mathrm{an}}$ et les \uncertain{appliquer} aux pts de $\mathbb{P}^1(\mathbb{C})$ pour trouver sections horizontales de $P$ ; la section horizontale qui passe par \struck{\ill{}} $p\in P_{x_0}$ s'obtient en choisissant un \struck{\ill{}} élément $f_{x_0}$ de $\mathcal{T}^{\mathrm{an}}_{x_0}$ ) d'où \struck{\ill{}} pt $p_0$ de $\mathbb{P}^1(\mathbb{C})_{X^{\mathrm{an}}}$ [\uncertain{tel que} $f_{x_0}$ : $\mathbb{P}^1(\mathbb{C})$ \struck{\ill{}} $\xrightarrow{\ \sim\ }P_x(\ni p)$ ], et définissant l'iso
\[
\text{pr.}\;x\mapsto\widetilde{f}_{x_0}(x),\qquad\text{où}\ \widetilde{f}_{x_0}\ \text{est la section}
\]
horizontale (de $P^{\mathrm{an}}$) passant par $f_{x_0}$. \uncertain{Or}

$P$ s'interprète comme $\mathcal{T}/B$ ($B\subset\mathrm{GP}(1)_{\mathbb{C}}$, Borel). Les sections horizontales \add{de $P^{\mathrm{an}}$} sont les sections horizontales de $\mathcal{T}^{\mathrm{an}}$ mod action de $B^{\mathrm{an}}$ — i.e.\ sections de $\mathcal{U}_x\subset\underline{\mathrm{Hol}}(X^{\mathrm{an}},\mathbb{P}^{1\,\mathrm{an}})$, mod action de $B^{\mathrm{an}}$, i.e.\ sections de $\mathcal{U}/B^{\mathrm{an}}$. Interprétation du $P$ disposant de $\sigma$ comme \uncertain{fournissant} une connexion \add{$D$} sur $\underline{\omega}$ \add{$D:\underline{\omega}\to\underline{\omega}^2$}, (on doit prendre les connexions $D$ telles que $c_D=c_0$ (la connexion projective donnée). Celles \add{(\uncertain{conj.} horizontales)} sont les sections de $\mathcal{T}/B_u$, et quel est
\note{la phrase se poursuit p. 111.}

\page{111}
\note{suite de la phrase de la p. 110.}
le fibré $\mathcal{T}/B_u$ \add{sur $P$} ? C'est le fibré \struck{\ill{}} \add{principal associé} : à $\Omega^1_{P/X}$ (ou encore des « éléments de contact projectifs d'ordre $1$ ») \add{\uncertain{restreint}} à une section en dessus de $\sigma$ est \struck{donc} \struck{une trivialisation de $\omega$} un couple d'un quotient inversible $L$ de $P^1_{X/S}(\underline{\omega})$ (qui donne $\sigma$) — i.e.\ une \struck{\ill{}} \add{\ill{}} de $\mathcal{U}$.
\[
D:\underline{\omega}\to L
\]
\struck{\ill{}} \add{d'ordre $1$} engendrant $L$, plus
\note{le dernier membre est souligné ; la ligne suivante est en partie biffée.}
une trivialisation $\varphi$ de $\mathcal{N}\otimes L^{-1}$, où $\mathcal{N}=\operatorname{Ker}(P^1_{X/S}(\underline{\omega})\to L)$, \struck{\ill{} $\simeq L^{-2}\otimes\underline{\omega}^3$ donc $\varphi:L^2\simeq\underline{\omega}^3$} — \marginal{i.e.\ $L$ est une « racine carrée de $\underline{\omega}^3$ »}
\struck{ou $\mathcal{N}\otimes L\simeq\det P^1_{X/S}(\underline{\omega})\simeq\underline{\omega}^3$ donc $D:\underline{\omega}\to\underline{\omega}^2$ engendrant $\underline{\mathcal{O}}_X$}
\note{passage biffé, lecture partielle.}

Le cas où $\sigma$ disjoint de $\Gamma$ \add{\uncertain{cf.}} correspond aux couples $(D,\varphi)$ d'une connexion \add{\uncertain{sur} $\underline{\omega}$} $D:\underline{\omega}\to\underline{\omega}^2$, et d'une trivialisation $\varphi$ de $\underline{\omega}^2$, $\varphi\in\Gamma(\underline{\omega})$. \struck{\uncertain{Comment} \ill{}} Condition d'horizontalité ? \struck{\ill{} horizontale \ill{} \ill{}} \add{\ill{} autre trivialisation} \add{\uncertain{triviale} la connexion $D$ sur $\omega$} \struck{\ill{}} i.e.\ $D\varphi=0$.

Tendance de la connexion projective $c_0$, quel est le faisceau $\mathcal{U}\subset\underline{\mathrm{Apphol}}(X^{\mathrm{an}},\mathbb{C}^{\mathrm{an}})$ ? C'est le faisceau des \add{germes} \struck{\ill{}} \add{\uncertain{étales}} \uncertain{d'applications} $t:X^{\mathrm{an}}\to\mathbb{P}^1_{\mathbb{C}}(\mathbb{C})$ telles que $t^*(\mathcal{P})=c_0$, où $c_0$ est la connexion projective (et unique !) sur $\mathbb{P}^1(\mathbb{C})$.

$\underline{\text{NB}}$ \ill{} $t^*(c)=c_{D_t}$, \struck{\ill{}} \uncertain{question} est donc $c_{D_t}=c_0$. \struck{\ill{}} Si $\xi:X'\to\mathbb{E}^1(\mathbb{C})\subset\mathbb{P}^1(\mathbb{C})$, est stable, on peut considérer la connexion \add{sur $\underline{\omega}$} $\xi^*(D)$ \struck{\ill{}}, où $D$ est la connexion canonique sur $\Omega^1_{\mathbb{P}^1(\mathbb{C})/\mathbb{C}}$ (la connexion $D_T$, où $T$ est la coordonnée…). C'est aussi $D_\xi$. On peut alors, \struck{\ill{}} pour une connexion \struck{\ill{}} $D:\underline{\omega}\to\underline{\omega}^2$ donnée, \struck{\ill{}} \add{$D$} \uncertain{regarder} les $t$ telles que $D=D_t$. \struck{\ill{}} Soit $t$ une uniformisante, telle \ill{} \ill{} donc
\[
D=D_t+\varpi\quad(\varpi\in\Gamma(\underline{\omega})),\ \text{et l'équation sur }s\mapsto D_s=D_t+\varpi,\ \text{on a vu que}
\]
\[
D_s=D_t+s'\,d\frac{1}{s'}\qquad s'=\frac{ds}{dt},\qquad s'\,d\frac{1}{s'}=-\frac{ds'}{s'}
\]
Donc l'équation est
\[
-\frac{ds'}{s'}=\varpi
\]
\note{sous $s'\,d\frac{1}{s'}$, l'auteur a d'abord écrit $-\frac{ds'}{s'^2}$, surchargé.}

\page{112}
En coordonnée, on pose $\varpi=\lambda(t)\,dt$ $\Rightarrow$
\[
\frac{s''}{s'}=-\lambda(t)\qquad\text{i.e.}\qquad(\log s')'=-\lambda(t)
\]
\[
\text{i.e.}\quad \log s'=-\Lambda(t)+\mathrm{c}^{\mathrm{te}}\qquad(\Lambda(t)\ \text{une primitive de}\ \lambda(t))
\]
\[
s'=\mathrm{c}^{\mathrm{te}}\exp\Lambda(t)\qquad s=\mathrm{c}^{\mathrm{te}}\int\exp-\Lambda(t)\,dt+\mathrm{c}^{\mathrm{te}}.
\]
\note{sous la première constante de $s'$ : « $\neq0$ ». Les signes de $\Lambda$ dans les deux formules ne concordent pas sur la page.}

Donc les solutions existent (\uncertain{analytiques}) et forment un torseur (sous le groupe affine de $\mathrm{Aut}(\mathbb{E}^1(\mathbb{C}))=(\mathrm{Aut}\,\mathbb{E}^1_{\mathbb{C}})^{\mathrm{an}}$ : c'est le \struck{\ill{}} \uncertain{\ill{}} \ill{}) transformant du torseur \struck{\ill{}} sous $B(\mathbb{C})^{-1}$ \struck{\ill{}} d'images inverses dans $\mathcal{T}$ des sections \uncertain{donnant} $\sigma'$ de $P$. Le torseur quotient de celui-ci par $B_u^{\mathrm{an}}\simeq\mathbb{G}_a^{\mathrm{an}}$ est le suivant : un section au-dessus de $\sigma'$ définie par $D$ est \struck{\ill{}} fournie de l'une des $1$-formes $\varphi$ de la forme $d\xi$, où $D_\xi=D$, i.e.\ des $\exists$ formes $\varphi$ telles que $\underline{D\varphi=0}$, $\varphi$ ne s'annulant pas.
\note{dans $D\varphi=0$, le $0$ est écrit sur un $\varphi$ biffé.}

\marginal{dans la marge gauche, en face des deux lignes suivantes : « \uncertain{analytique} » et « $\mathrm{NB}$ »}
Sections \add{$\simeq\mathcal{T}/B$,} \struck{\uncertain{analytiques}} de $P$ horizontales disjointes de $\sigma$ : connexions \add{sur $\underline{\omega}$} telles que $c_D=c_0$ \struck{\ill{}} $\simeq$ \uncertain{régulières} de $\mathcal{T}/B_u$ horizontales, \uncertain{ou} encore

\marginal{dans la marge gauche, le long de l'accolade : « \uncertain{analytique} — transcendante »}
Sections \emph{analytiques} \uncertain{de} sections $\sigma'$ disjointes de $\sigma$ : couples \struck{\ill{}} $(D,\varphi)$, $D$ connexion \add{\emph{analyt.} sur $\underline{\omega}$} telle que $c_D=c_0$, $\varphi\in\Gamma^{\mathrm{an}}(\underline{\omega})$ telle que $D\varphi=0$ \add{telle que $\varphi$ \uncertain{s'annule pas}}.

Sections analytiques de $\mathcal{T}$ horizontales : \uncertain{uniformisantes} \struck{$t$ telles que} $\xi:X\to\mathbb{E}^1_{\mathbb{C}}$ telles que $s^*(\mathcal{C})=c_0$, i.e.\ $c_{D_t}=c_0$.
\[
t\longmapsto(D_t,dt)\longmapsto D_t\qquad\text{les applications canoniques}
\]
\[
\text{correspondent :}\qquad \mathcal{T}\to\mathcal{T}/B_u\to\mathcal{T}/B
\]

\page{113}
\note{feuille double écrite en deux colonnes ; la colonne de gauche est transcrite d'abord.}
Équations explicites, en termes d'un paramètre local $t:X\to\mathbb{E}^1_{\mathbb{C}}$ de sorte que
\[
c_0=c_{D_t}+\varpi_2\qquad\varpi_2\in\Gamma(\underline{\omega}^2),\quad\varpi_2=\mu(t)\,dt^2
\]

a) Sections horizontales $\sigma'$ de $P=\mathcal{T}/B$ \add{disjointes de $\sigma$}, correspondant aux conn.\ $D$ sur $\underline{\omega}$ telles que $c_D=c_0$, i.e.\ (posant $D=D_t+\omega$) aux $1$-formes $\omega$ telles que $c_{D_t+\omega}=c_0$ i.e.\ (puisque $c_{D_t+\omega}=c_{D_t}+2D_t\omega+\omega^2$, $c_0=c_{D_t}+\varpi_2$)
\[
2D_t\omega+\omega^2=\varpi_2
\]
ou encore, posant $\omega=\lambda(t)\,dt$
\[
(1)\qquad\boxed{2\lambda'(t)+\lambda(t)^2=\mu(t)}\qquad\Bigl(\text{où }\lambda'=\frac{d}{dt}\lambda\Bigr)
\]

b) Sections horizontales de $\mathcal{T}/B_u$, au-dessus d'une section $\sigma'$ de $\mathcal{T}/B$ disjointe de $\sigma$. Couples $(D,\varphi)$, avec $D$ comme dessus, $\varphi\in\Gamma(\underline{\omega})^*$, $D\varphi=0$. Posant $\varphi=\alpha(t)\,dt$, on a
\[
D\varphi=D_t\varphi+\omega\varphi=(\alpha'(t)+\alpha(t)\lambda(t))\,dt^2
\]
donc
\[
(2)\qquad\boxed{\alpha'(t)+\alpha(t)\lambda(t)=0}\qquad\Bigl(\text{où }\alpha'=\frac{d\alpha}{dt}\Bigr)\qquad\text{i.e.}\quad\frac{\alpha'}{\alpha}=-\lambda
\]
\[
\text{i.e.}\quad\frac{\alpha'(t)}{\alpha(t)}=-\lambda(t),\quad\text{ou}\quad\log\alpha(t)=-\int\lambda(t)\,dt,
\]
\[
\text{ou}\quad\alpha(t)=c\exp-\int\lambda(t)\,dt\qquad c\ \mathrm{c}^{\mathrm{te}}\neq0
\]

d) Sections horizontales de $\mathcal{T}$, \struck{\ill{}} correspondant à une section $\sigma'$ de $\mathcal{T}/B$ disjointe de $\sigma$ : $\xi:X^{\mathrm{an}}\to\mathbb{E}^1(\mathbb{C})$ telle que $c_{D_\xi}=c_0$, i.e.
\[
c_{D_t}\;\text{\struck{$-\frac{ds'}{s'}\,dt$}}=c_{D_t}+\varpi_2
\]
i.e.\ (\uncertain{déjà})
\[
2D_t\Bigl(-\frac{ds'}{s'}\Bigr)+\Bigl(-\frac{ds'}{s'}\Bigr)^2=\varpi_2\qquad\Bigl(s'=\frac{ds}{dt}\Bigr)
\]
qui s'explicite aussi en (4) : $-\dfrac{ds'}{s'}=-\dfrac{s''}{s'}=\lambda(t)\,dt$.
\note{le passage d) est encadré d'un trait ; devant $-\frac{ds'}{s'}$, un premier essai biffé.}

c) Sections horizontales de $\mathcal{T}/B_u$ \struck{\ill{}} au-dessus de $\sigma'$ définie par $D$ ($=D_t+\omega=D_t+\lambda(t)\,dt$, où $\lambda(t)$ solution de (1) i.e.\ $\omega$ solution de $\alpha$) : les $s:X^{\mathrm{an}}\to\mathbb{E}^1_{\mathbb{C}}=\mathbb{C}$ telles que $D_s=D$, i.e.\ (posant $s=s(t)$) \uncertain{avec} $\frac{ds'}{s'}=-\omega$ ($s'=\frac{ds}{dt}$)
\[
(3)\qquad\frac{s''}{s'}=-\lambda(t)\qquad s'=\frac{ds}{dt},\quad s''=\frac{ds'}{dt}
\]
Pour passer de b') à b) : poser $\varphi=ds=s'\,dt$, i.e.\ $\boxed{\alpha=s'}$.
\note{les lettres a) à e) ne suivent pas exactement l'ordre de la p. 112 ; la colonne de gauche les donne dans l'ordre a), b), d), c).}

\[
2\Bigl(-\frac{s''}{s'}\Bigr)'+\Bigl(-\frac{s''}{s'}\Bigr)^2=\mu(t)
\]
\[
\text{i.e.}\quad -2\,\frac{s'''s'-s''^2}{s'^2}+\frac{s''^2}{s'^2}=\mu
\]
\[
\frac{3s''^2-2s's'''}{s'^2}=-2\,\frac{s'''}{s'}+3\Bigl(\frac{s''}{s'}\Bigr)^2
\]
\[
(4)\qquad\boxed{G_t(s)\overset{\mathrm{dfn}}{=}-2\,\frac{s'''}{s'}+3\Bigl(\frac{s''}{s'}\Bigr)^2=\mu(t)}
\]
\note{au-dessus de l'encadré, un essai biffé : « $s''=-2s'\mu(t)$ », lecture incertaine.}
où $s'=\frac{ds}{dt}$, $s''=\frac{ds'}{dt}$, $s'''=\frac{ds''}{dt}$.

Pour passer de d) à c) : $s\longmapsto(ds,\ \lambda=-\frac{s''}{s'})$

\qquad de d) à b) : $s\longmapsto(\alpha=s',\ \lambda=-\frac{s''}{s'})$

\qquad de d) à a) : $s\longmapsto\lambda=-\frac{s''}{s'}$

$\underline{\text{NB}}$ La propriété essentielle de l'expression différentielle $G_t(s)=-2\frac{s'''}{s'}+3\bigl(\frac{s''}{s'}\bigr)^2$ \struck{$(=(D_s-D_t)/dt^2$} ($=(c_{D_s}-c_{D_t})/dt^2$) est son invariance par le groupe homographique : si
\[
\text{\struck{$\overline{s}=as$}}\qquad\overline{s}=\frac{as+b}{cs+d}\qquad a,b,c,d\ \text{constantes},
\]
alors $G_t(\overline{s})=G_t(s)$ … De plus, si $\tau$ est une autre uniformisante, on a
\[
G_\tau(s)=G_t(s)\,\tau'^2
\]
\note{le facteur $\tau'^2$ suit un signe surchargé illisible ; la formule n'est pas corrigée ici.}

$\underline{\text{Th}}$ Soit $X$ courbe alg.\ complète sur $\mathbb{C}$, $t$ une uniformisante sur un ouvert $\mathcal{U}$ de $X$, i.e.\ morphisme étale $t:\mathcal{U}\to\mathbb{E}^1_{\mathbb{C}}$. Il existe alors une unique section $\mu$ de $\underline{\omega}_{\mathcal{U}}$ \uncertain{ayant} la propriété suivante : Considérons l'équation \struck{\ill{}} différentielle sur $\mathcal{U}^{\mathrm{an}}$, en la fonction analytique inconnue $s:\mathcal{U}\to\mathbb{E}^1_{\mathbb{C}}$
\[
G_t(s)\overset{\mathrm{dfn}}{=}-2\,\frac{s'''}{s'}+3\Bigl(\frac{s''}{s'}\Bigr)^2=\mu(t)
\]
Alors les solutions de cette équation, sur \uncertain{tout}
\note{la phrase se poursuit p. 114.}

\page{114}
\note{suite de l'énoncé de la p. 113.}
ouvert \add{$V$} de $\mathcal{U}^{\mathrm{an}}$, sont des \uncertain{restrictions}, à $V$, de fonctions \uncertain{uniformisantes} \add{\uncertain{\ill{}}} sur $X^{\mathrm{an}}$ (et \struck{\ill{}} \add{\ill{}} cor.\ exactement toutes les fonctions uniformisantes \uncertain{\ill{}}).

\emph{Corollaire} Si $X$ est \uncertain{\ill{}} complète, \ill{} le \ill{} \ill{} \uncertain{seul} $\mathcal{D}(\mu)$ \uncertain{tel} que $\mu(t)$ est \uncertain{restriction} d'une section \emph{analytique} de $\underline{\omega}_{\mathcal{U}}$. \struck{\ill{} \ill{}} \add{Est-elle} algébrique ? Premier cas douteux : $X=\mathbb{P}^1_{\mathbb{C}}-$ trois pts, où $X$ = courbe elliptique moins un point.

\emph{Autre question} Quand $X$ (supposée complète disons), varie avec un paramètre \emph{algébrique} parcourant une variété algébrique $S$, int.\ \ill{} $X\to S$, \uncertain{est-ce} \uncertain{donc} $\mu$ des fonctions analytiques sur \ill{} de $X$, algébriques sur les fibres. Sont-elles algébriques ? Ou bien (\uncertain{si} \uncertain{non}), est-elle une signification \emph{transcendante} algébrique, indépendante de la top.\ des corps \ill{} $S$ ou $\mathcal{C}$ — dans laquelle un \ill{} plus fin que $S$ est une \uncertain{relation} \ill{} de \ill{} \uncertain{celle} ?).

\page{115}
\note{mise au net, d'une écriture soignée, de la théorie de la dérivée schwarzienne.}
Soit $\mathcal{X}$ une « courbe analytique » i.e.\ espace \uncertain{\ill{}} analytique régulier partout de dim.\ $1$. Soit
\[
t:\mathcal{X}\to\mathbb{C}
\]
une « uniformisante » i.e.\ une application analytique étale de $X$ dans $\mathbb{C}=\mathbb{E}^1(\mathbb{C})$, i.e.\ une section $t$ de $\underline{\mathcal{O}}_X$ telle que $dt\in\Gamma(X,\underline{\Omega}^1_{X/\mathbb{C}})$ soit partout $\neq0$ (i.e.\ soit une \emph{base} de $\underline{\Omega}^1_{X/\mathbb{C}}$). Soit $s\in\Gamma(\mathcal{X},\underline{\mathcal{O}}_{\mathcal{X}})$ — \struck{\ill{}} une autre uniformisante, et posons
\[
G(t,s)=\Bigl(-2\,\frac{s'''}{s'}+3\Bigl(\frac{s''}{s'}\Bigr)^2\Bigr)dt^2\in\Gamma(\mathcal{X},\underline{\omega}^{\otimes2})
\]
où $s'=\frac{ds}{dt}$, $s''=\frac{ds'}{dt}$, $s'''=\frac{ds''}{dt}$. On \struck{\ill{}} vérifie facilement, si $a\in\mathbb{C}^*$, $b\in\mathbb{C}$ \struck{que si $a,b$ sont des constantes, $a\neq0$ et}
\[
(1)\qquad
\left\{
\begin{aligned}
&G(t,s+b)=G(t,s) &&\text{a)}\\
&G(t,as)=G(t,s) &&\text{b)}\\
&G\bigl(t,\tfrac1s\bigr)=G(t,s) &&\text{c)}
\end{aligned}
\right.
\]
où, dans la dernière équation, on suppose que $\frac1s$ est aussi uniformisant, i.e.\ $s$ ne prend pas de zéros sur $\mathcal{X}$. On vérifie de plus que, si $t,s,r$ sont trois uniformisantes,
\[
(2)\qquad G(t,s)+G(s,r)=G(t,r)
\]
\uncertain{Enfin}, \add{(3)} la formation de $G(t,s)$ commute à la restriction à un ouvert. \uncertain{En}

De (1) et (3), résulte qu'on peut prolonger l'expression $G(t,s)$ au cas où $s$, au lieu d'être une « uniformisante » i.e.\ une application étale $\mathcal{X}\to\mathbb{E}^1(\mathbb{C})$, est une application étale
\[
\mathcal{X}\to\mathbb{P}^1(\mathbb{C})
\]
i.e.\ une fonction méromorphe n'ayant que des pôles simples, et uniformisante en dehors de ceux-ci — (dans ce cas, \add{\uncertain{on vérifie que} dans} \uncertain{voisinage} des $s$ où $s$ a un pôle $G(t,s)=G(t,\frac1s)$), et dans (1) \struck{\ill{}} \add{\uncertain{reste} vraie avec $s$ méromorphe} \struck{\ill{}} \add{\ill{}} reste valable, \struck{et} (plus généralement)
\[
G\Bigl(t,\frac{as+b}{cs+d}\Bigr)=G(t,s)\qquad a,b,c,d\in\mathbb{C},\quad ad-bc\in\mathbb{C}^*
\]

\page{116}
D'autre part, faisant $t=r$ dans (2), on trouve
\[
G(t,t)=0\qquad(t\ \text{uniformisante})
\]
et faisant $r=t$, on trouve
\[
G(t,s)+G(s,t)=0,\quad\text{i.e.}\quad G(t,s)=-G(s,t)\qquad(s,t\ \text{uniformisantes})
\]
\struck{Donc on peut prolonger l'expression $G(t,s)$, définie \ill{} au cas où $t:\mathcal{X}\to\mathbb{E}^1(\mathbb{C})$, $s:\mathcal{X}\to\mathbb{P}^1_{\mathbb{C}}(\mathbb{C})$, au cas aussi où \ill{} \ill{} des pôles, \ill{}}
\note{passage encadré et biffé ; dans la marge gauche, en oblique, une phrase biffée elle aussi : « au cas où $t$ \ill{} c'est \ill{} \ill{} $\mathcal{X}\to\mathbb{P}^1(\mathbb{C})$ ».}
d'où
\[
(1\,\text{bis})\qquad
\left\{
\begin{aligned}
&G(at,s)=G(t,s)\\
&G(t+b,s)=G(t,s)\\
&G\bigl(\tfrac1t,s\bigr)=G(t,s)\qquad\bigl(t,\tfrac1t\ \text{unif}^{\text{tes}}\bigr)
\end{aligned}
\right.
\]
ce qui permet de faire \uncertain{la} généralisation pour les $t$, comme pour les $s$, d'où finalement
\[
s,t:\mathcal{X}\to\mathbb{P}^1(\mathbb{C})\qquad\text{morphismes étales}
\]
\[
(4)\qquad
\left\{
\begin{aligned}
&G\Bigl(\frac{at+b}{ct+d},s\Bigr)=G(t,s)\qquad a,b,c,d\in\mathbb{C}\ \text{avec}\ ad-bc\in\mathbb{C}^*\\
&\text{de plus,}\quad G\Bigl(t,\frac{as+b}{cs+d}\Bigr)=G(t,s)\qquad\text{id.}
\end{aligned}
\right.
\]

\struck{De plus, on a} Soit $\mathcal{P}_X$ le faisceau des morphismes \add{analytiques} étales
\[
\mathcal{X}\to\mathbb{P}^1(\mathbb{C}),
\]
sur lequel opère le groupe
\[
H=\mathrm{Aut}_{\mathbb{C}}\mathbb{P}^1(\mathbb{C})\simeq\mathrm{GP}(1,\mathbb{C})=\Bigl\{\begin{pmatrix}a&b\\c&d\end{pmatrix}\Bigm|ad-bc\in\mathbb{C}^*\Bigr\}
\]
(mod.\ matrices diagonales). On a donc trouvé un morphisme de faisceaux
\[
\underline{G}:\mathcal{P}_X\times\mathcal{P}_X\to\underline{\omega}_X^{\otimes2}
\]
satisfaisant les deux conditions :
\begin{enumerate}
\item[I)] $G$ invariant par $H\times H$ i.e.
\[
G(g\cdot t,g'\cdot s)=G(t,s)\quad\text{si }g,g'\in H
\]
\item[II)] $G(t,s)+G(s,r)=G(t,r)$
\end{enumerate}
En particulier $G(t,t)=0$, $G(t,s)=-G(s,t)$.

En fait, le théorème des équations différentielles (\uncertain{donne}, pour $t$ fixé), donne aisément une réciproque : I), savoir
\note{la phrase se poursuit p. 117.}

\page{117}
\note{suite de la phrase de la p. 116.}
\begin{enumerate}
\item[I')] a) si $s,t\in\Gamma(\mathcal{U},\mathcal{P}_X)$, \add{$\mathcal{U}$ \uncertain{connexe}} \struck{\ill{}} et $G(t,s)=0$, alors $\exists$ \struck{\ill{}} $g\in H$ tel que $t=g\cdot s$.

b) Pour $t$ fixé, le morphisme
\[
s\longmapsto G(t,s)\qquad\mathcal{P}_X\to\underline{\omega}_X^{\otimes2}
\]
(qui induit un homom.\ injectif \struck{$\mathcal{P}_X/H$}
\[
H\backslash\mathcal{P}_X\hookrightarrow\underline{\omega}_X^{\otimes2}
\]
\uncertain{d'après} a)) est \emph{surjectif} (i.e.
\[
H\backslash\mathcal{P}_{\mathcal{X}}\xrightarrow{\ \sim\ }\underline{\omega}_X^{\otimes2}\qquad s\longmapsto G(t,s)=\Bigl(-2\,\frac{s'''}{s'}+3\Bigl(\frac{s''}{s'}\Bigr)^2\Bigr)dt^2\ \bigr)
\]
\end{enumerate}

D'autre part, grâce à II, on peut donner une forme plus simple et intrinsèque à l'expression $G(t,s)$.

\emph{Proposition} (\uncertain{presque} formelle à partir de (II)) On peut, de façon unique : trouver un torseur $\mathrm{Conn}_{\mathcal{X}}$ sous $\underline{\omega}_{\mathcal{X}}^{\otimes2}$, et un morphisme de faisceaux
\[
t\longmapsto c_t:\ \mathcal{P}_X\xrightarrow{\ c^X\ }\mathrm{Conn}_{\mathcal{X}}
\]
tels que l'on ait, pour $s,t\in\Gamma(\mathcal{U},\mathcal{P}_X)$
\[
\boxed{G(t,s)=c_s-c_t}
\]

Le faisceau $\mathrm{Conn}_{\mathcal{X}}$ est appelé le faisceau des \emph{connexions projectives} (analytiques complexes) sur $\mathcal{X}$. Une section de ce dernier est appelée connexion projective (analytique) de $\mathcal{X}$. \struck{\ill{}} Le morphisme $c^X$ est invariant par $H$ et \uncertain{définit} un iso.\ canon.\ :
\[
c^X:H\backslash\mathcal{P}_{\mathcal{X}}\xrightarrow{\ \sim\ }\mathrm{Conn}_{\mathcal{X}}
\]
\marginal{une flèche relie « $\mathcal{P}_X$ » au mot « une uniformisante étale ».}
Une uniformisante étale
\[
f:\mathcal{X}'\to\mathcal{X}
\]
définit des iso
\[
\mathcal{P}_{\mathcal{X}'}\simeq f^{-1}(\mathcal{P}_{\mathcal{X}})\qquad\mathrm{Conn}_{\mathcal{X}'}\simeq f^{-1}(\mathrm{Conn}_{\mathcal{X}})
\]
et on trouve un \struck{\ill{}} diagramme commutatif

\begin{tikzcd}
\mathcal{P}_{\mathcal{X}'} \arrow[r, "c^{\mathcal{X}'}"] \arrow[d, "\wr"'] & \mathrm{Conn}_{\mathcal{X}'} \arrow[d, "\wr"] \\
f^{-1}(\mathcal{P}_{\mathcal{X}}) \arrow[r, "c^{\mathcal{X}}"'] & f^{-1}(\mathrm{Conn}_{\mathcal{X}})
\end{tikzcd}

\note{sur la flèche du haut, un premier label biffé.}

\page{118}
Notons que si $\mathcal{X}_0=\mathbb{P}^1(\mathbb{C})$ \struck{il y a le} faisceau $H\backslash\mathcal{P}_{\mathcal{X}_0}$ est le faisceau final ; \struck{et que} il y a une conn.\ projective can.\ \add{$\mathcal{C}$} sur $\mathcal{X}_0$, savoir $c_t$ pour $t=\mathrm{id}_{\mathcal{X}_0}$ ! C'est la seule conn.\ projective sur $\mathcal{X}_0$, car $H^0(\mathcal{X}_0,\underline{\omega}^{\otimes2}_{\mathcal{X}_0})=0$ ($\underline{\omega}$ est de degré $-2$, donc $\underline{\omega}^{\otimes2}$ de degré $-4$ !).

Revenant au cas d'une $\mathcal{X}$ générale, on voit donc que
\[
c_t=t^*(\mathcal{C})=\text{image inverse par }t\text{ de la conn.\ projective can.\ }\mathcal{C}
\]
(et unique conn.\ proj.) sur $\mathbb{P}^1(\mathbb{C})$.

\emph{Extension à la géom.\ alg.} \quad \uncertain{Soit} maintenant $X$ schéma formel. \uncertain{Lisse} \add{ou plus} de dimension relative $1$ sur un schéma $S$. Si \struck{$t,s$} \ill{}
\[
t,s:X\to\mathbb{E}^1_S
\]
\uncertain{étales} sont \add{\uncertain{\ill{}}} étales, on définit comme plus haut $G(t,s)\in\Gamma(X,(\underline{\Omega}^1_{X/S})^{\otimes2})$ \add{$\simeq\underline{\omega}^{\otimes2}$}, et on vérifie bien \uncertain{vite} (1) et (2) pour $s=$ \uncertain{\ill{}} \uncertain{\ill{}} que dans le cas analytique, — qui permet d'étendre $G(t,s)$ au cas de (uniformisantes \emph{étales})
\[
t,s:X\to\mathbb{P}^1_S
\]
\struck{\ill{}}, de façon à satisfaire encore \struck{$t,s:X\to$}
\[
\text{(\uncertain{groupe} }H_S=\underline{\mathrm{Aut}}_S(\mathbb{P}^1_S))
\]
\[
G(gt,g's)=G(t,s)\qquad(g,g'\in\Gamma(S,H))
\]
et : définissant ainsi, si $\mathcal{P}_{X/S}$ est le faisceau \add{(\uncertain{étale})} des uniformisantes étales $X\to\mathbb{P}^1_S$, \struck{des \ill{}}
\[
G:\mathcal{P}_{X/S}\times\mathcal{P}_{X/S}\to\underline{\omega}^{\otimes2}_{X/S}
\]
\marginal{dans la marge gauche, en oblique : « en car.\ \uncertain{nulle} \ill{} ou \ill{} \ill{} \ill{} \ill{} \ill{} \ill{} $\geqslant3$, i.e.\ \uncertain{que} $P^3_{X/S}(\underline{\omega})$ \ill{} \ill{} $S$ \ill{} $\geqslant3$ »}

\page{119}
\emph{invariante} sous l'action de $f^{-1}(H_S)$ i.e.\ telle que
\[
G(g\cdot s,g'\cdot t)=G(s,t)\quad\text{si}\quad
\left\{
\begin{aligned}
&s,t\in\Gamma(\mathcal{U},\mathcal{P}_X)\\
&g,g'\in\Gamma(V,H_V)
\end{aligned}
\right.
\]
où

\begin{tikzcd}
U \arrow[r, "\text{ét}"] \arrow[d] & X \arrow[d, "f"] \\
V \arrow[r, "\text{ét}"'] & S
\end{tikzcd}

\note{le carré est marqué commutatif par une flèche courbe entre les deux verticales.}
et satisfaisant
\[
G(s,t)+G(t,r)=G(s,r).
\]

Cela permet donc de définir le faisceau (étale) $\mathrm{Conn}_{X/S}$ des connexions projectives sur $X$ rel.\ à $S$, qui est un torseur sous $\underline{\omega}^{\otimes2}_{X/S}$, et d'interpréter $G(t,s)$, à l'aide d'une uniformisante canonique
\[
c^X:\mathcal{P}_{X/S}\to\mathrm{Conn}_{X/S}
\]
par
\[
G(t,s)=c^X(s)-c^X(t)
\]

Enfin, par construction, la formation de $\mathcal{P}_{X/S}$, $\mathrm{Conn}_{X/S}$ et $c^X$ commute à la loc.\ étale. Cela permet d'interpréter comme
\[
c^X(t)=t^*(\mathcal{C})
\]
où $\mathcal{C}$ est l'unique connexion sur $\mathbb{P}^1_S$.

Plus généralement, pour $\mathcal{X}_0$ un fibré en droites projectives $\mathcal{X}_0$ sur $S$, il y a une unique conn.\ proj.\ $c_{\mathcal{X}_0}$ sur $\mathcal{X}_0$, et si $t,s:X\to\mathcal{X}_0$ sont étales, on peut définir
\[
G(t,s)=c(s)-c(t)=s^*(c_{\mathcal{X}_0})-t^*(c_{\mathcal{X}_0})
\]
satisfaisant aux conditions analogues : celles du cas $\mathcal{X}_0=\mathbb{P}^1_S$, que le lecteur (sic) écrira…
\note{une flèche relie la formule $G(t,s)=\dots$ à « conditions analogues ».}

\page{120}
\note{feuille double écrite en deux colonnes ; la colonne de gauche est transcrite d'abord.}
Si $S$ est de car.\ $0$, alors si $X$ est \add{(lisse$/S$ ?)} à fibres géom.\ connexes non vides, on peut montrer que pour $s,t:X\to\mathcal{X}_0$ (fibré en droites projectives sur $S$) étales, tels que $G(t,s)=0$ i.e.\ $t^*(c_{\mathcal{X}_0})=s^*(c_{\mathcal{X}_0})$, il existe
\[
g\in\mathrm{Aut}(\mathcal{X}_0)\quad\text{tel que}\quad s=g\circ t
\]
\marginal{dans la marge gauche, contre ces lignes, barré d'un trait oblique : « \uncertain{unicité} ! »}
Donc, pour $t$ fixé, $s$ variable, le morphisme
\[
f^{-1}(H_S)\backslash\mathcal{P}_X\to\underline{\omega}^{\otimes2}_{X/S}\qquad s\longmapsto G(t,s)
\]
est \emph{injectif}, ou encore (\uncertain{indépendamment}) le morphisme
\[
{}_H\backslash\mathcal{P}_X\xrightarrow{\ c^X\ }\mathrm{Conn}_{X/S}
\]
défini par passage au quotient par $c^X$, est \emph{injectif}. \struck{\ill{}} Mais bien sûr ces deux morphismes ne sont plus \struck{bijectifs} surjectifs. Une conn.\ proj.\ qui est \struck{contenue dans} dans l'image de ${}_H\backslash\mathcal{P}_X$ est dite « intégrable par fonctions algébriques ».

Si $S$ est de car.\ $p>0$, \uncertain{alors} l'injectivité n'est plus valable non plus. Introduisons $X^{(p/S)}$ et
\[
X\xrightarrow{\ F_X\ }X^{(p/S)}\qquad(\text{Frobenius})
\]
\note{les deux schémas sont reliés à $S$ par deux traits obliques, formant un triangle.}
On trouve que \struck{\ill{} \ill{} \ill{}} le groupe des $H(X^{(p/S)})$, qui est égal \struck{de façon évidente}, sur $\mathcal{P}_X$, $\Gamma(X,\mathcal{P}_X)=\mathrm{Et}_S(X,\mathbb{P}^1_S)$, laisse invariant l'expression $G(t,s)$, et le mieux qu'on puisse espérer, c'est d'avoir
\[
F_X^{-1}(H_{X^{(p/S)}})\backslash\mathcal{P}_X\hookrightarrow\mathrm{Conn}_{X/S}
\]
\marginal{dans la marge gauche, en oblique : « si $2$ et $3$ inversibles, \uncertain{\ill{} \ill{}} ! cf.\ plus bas »}

Introduisons (\uncertain{pour} $s,t$ \struck{\ill{}} uniformisantes) l'expression
\[
D^{(n)}_tf=\frac{1}{n!}\,\frac{d^nf}{dt^n}
\]
(qui a un sens sans restriction de car.\ \uncertain{nulle}…, dès que $X/S$ est diff.\ lisse de \struck{\ill{}} rel.\ dim.\ $1$) on trouve alors (pour toutes uniformisantes)
\[
G(t,s)=12\Bigl(-\frac{D^{(3)}_ts}{D^{(1)}_ts}+\Bigl(\frac{D^{(2)}_ts}{D^{(1)}_ts}\Bigr)^2\Bigr)=12\,G_0(t,s)
\]
où
\[
G_0(t,s)=-\frac{D^{(3)}_ts}{D^{(1)}_ts}+\Bigl(\frac{D^{(2)}_ts}{D^{(1)}_ts}\Bigr)^2
\]
Cela montre que si $12\cdot1_S$ \struck{$\neq0$}, dans toutes les conn.\ projectives de la forme $c_t$ sont égales, donc il y a une conn.\ proj.\ canonique sur $X$ rel.\ $S$, à laquelle on peut \uncertain{infirmer} toutes les autres. Cela montre aussi (sans restriction de car.\ telle $12\cdot1_S=0$ !) que le torseur \struck{\ill{}} $\mathrm{Conn}_{X/S}$ est déduit moyennant le changement de groupe \uncertain{déf.}
\[
\underline{\omega}^{\otimes2}_{X/S}\xrightarrow{\ 12\,\mathrm{id}\ }\underline{\omega}^{\otimes2}_{X/S},
\]
d'un torseur plus fin $\mathrm{Conn}_{0,X/S}$, sous $\underline{\omega}^{\otimes2}_{X/S}$, avec un morphisme
\[
\mathcal{P}_{X/S}\xrightarrow{\ c^X_0\ }\mathrm{Conn}_{0,X/S}
\]
tel que
\[
c^X_0(s)-c^X_0(t)=G_0(t,s)
\]
Bien entendu, les \uncertain{assertions} plus haut sur \add{l'expression} $G(t,s)$ et les constructions
\note{la phrase se poursuit au-delà du lot.}

\end{document}
